% archive/manual-glossary/app-a-glossary.tex
% Archived manual glossary snapshot.
% Historical snapshot from the old autogenerated glossary pipeline.

\chapter{Глоссарий}
\label{app:glossary}

Термины приведены в алфавитном порядке: сначала латиница (A--Z),
затем кириллица (А--Я). Ссылка на главу указывает, где термин
раскрыт подробно.

%% ═══════════════════════════════════════════════════════════════
%%  ЛАТИНИЦА  A – Z
%% ═══════════════════════════════════════════════════════════════

\section*{A}

\textbf{AAC} (Application Authentication Cryptogram) —
криптограмма EMV-чипа, означающая \enquote{карта отклоняет
  транзакцию}; терминал не должен продолжать
операцию. См. главу \ref{ch:emv}.

\textbf{ACS} (Access Control Server) — сервер в домене
эмитента, принимающий решение об аутентификации держателя
в протоколе 3-D Secure. ACS оценивает риск транзакции и выбирает
между frictionless и challenge flow. См. главу \ref{ch:3ds}.

\textbf{AID} (Application Identifier) — уникальный
идентификатор платёжного приложения на EMV-чипе; состоит из RID
(5 байт, идентификатор карточной сети) и PIX (до 11 байт,
конкретный продукт). См. главу \ref{ch:emv}.

\textbf{AML} (Anti-Money Laundering) — совокупность мер,
процедур и систем, направленных на предотвращение легализации
(отмывания) доходов, полученных преступным
путём. См. главу \ref{ch:compliance-aml}.

\textbf{ARQC} (Authorization Request Cryptogram) —
криптограмма EMV-чипа, означающая \enquote{карта просит
  онлайн-авторизацию у эмитента}; самый распространённый исход
контактной EMV-транзакции. См. главу \ref{ch:emv}.

\textbf{ATC} (Application Transaction Counter) — счётчик
на EMV-чипе, инкрементируемый при каждой транзакции; используется
для обнаружения клонирования и replay-атак. См. главу \ref{ch:emv}.

\section*{B}

\textbf{BIN} (Bank Identification Number) — первые шесть
(исторически) или восемь (с 2022 года) цифр PAN,
идентифицирующие эмитента и карточную сеть. Формально
переименован в IIN по \iso{7812}, но термин BIN по-прежнему
широко используется в индустрии. См. главу \ref{ch:card-data}.

\textbf{Bitmap} — битовая карта (64 или 128 бит) в сообщении
\iso{8583}, указывающая, какие data elements присутствуют
в данном сообщении. См. главу \ref{ch:iso8583}.

\section*{C}

\textbf{CDE} (Cardholder Data Environment) — совокупность
людей, процессов и технологий, которые хранят, обрабатывают
или передают cardholder data или sensitive authentication data,
а также все системы и сегменты сети, напрямую подключённые к ним
или влияющие на их безопасность. См. главу \ref{ch:pci-dss}.

\textbf{Chargeback} — формальный диспут: принудительный возврат
средств по транзакции, инициированный эмитентом по заявлению
держателя карты. См. главу \ref{ch:disputes}.

\textbf{Clearing} (клиринг) — пакетная обработка транзакций
после авторизации: формирование клиринговых файлов,
сопоставление с авторизациями и расчёт
interchange. См. главу \ref{ch:tx-lifecycle}.

\textbf{CNP} (Card-Not-Present) — операция без физического
присутствия карты: интернет-покупка, заказ по телефону
или по почте. См. главу \ref{ch:card-data}.

\textbf{CVM} (Cardholder Verification Method) — способ
верификации держателя карты: PIN, подпись, биометрия
или без верификации; выбирается терминалом на основе конфигурации
чипа и правил сети. См. главу \ref{ch:emv}.

\textbf{CVV} (Card Verification Value) — код верификации,
записанный на магнитной полосе карты; вычисляется эмитентом
на основе PAN, срока действия и сервисного кода.
У \scheme{Mastercard} аналог называется
CVC. См. главу \ref{ch:card-data}.

\textbf{CVV2} / \textbf{CVC2} — трёхзначный (или четырёхзначный
у American Express) код, напечатанный на карте, предназначенный
для CNP-операций. Вычисляется иначе, чем CVV, и всегда
от него отличается. См. главу \ref{ch:card-data}.

\section*{D}

\textbf{Data Element} (DE) — индивидуальное поле в сообщении
\iso{8583} с фиксированным номером и назначением (например,
\de{2} — PAN, \de{4} — сумма, \de{39} — код
ответа). См. главу \ref{ch:iso8583}.

\textbf{DCC} (Dynamic Currency Conversion) — конвертация
валюты в точке продажи: мерчант предлагает держателю карты
оплатить в валюте карточного счёта вместо валюты
мерчанта. См. главу \ref{ch:multicurrency}.

\textbf{dCVV} / \textbf{CVC3} — динамический код верификации,
генерируемый EMV-чипом для каждой бесконтактной транзакции
с использованием сессионного ключа и unpredictable
number. См. главу \ref{ch:card-data}.

\textbf{DDA} (Dynamic Data Authentication) — метод офлайн-аутентификации
EMV-чипа, при котором карта формирует подпись с использованием
данных, уникальных для каждой
транзакции. См. главу \ref{ch:emv}.

\textbf{DPAN} (Device PAN, Digital PAN) — токен, заменяющий
настоящий PAN в платёжной транзакции; выпускается Token Service
Provider карточной сети, проходит проверку Луна и имеет формат
стандартного PAN, но принадлежит отдельному
BIN-диапазону. См. главу \ref{ch:tokenization}.

\textbf{DS} (Directory Server) — сервер в домене
совместимости 3-D Secure, управляемый карточной сетью;
маршрутизирует аутентификационные сообщения между 3DS Server
мерчанта и ACS эмитента, но не принимает решение
об аутентификации. См. главу \ref{ch:3ds}.

\textbf{Dual message} — архитектура обработки транзакций,
в которой авторизация и клиринг — два отдельных сообщения,
разнесённых во времени (MTI \mti{0100}/\mti{0110}); стандарт
для POS-покупок. См. главу \ref{ch:tx-lifecycle}.

\section*{E}

\textbf{EMV} — стандарт для микропроцессорных платёжных карт,
обеспечивающий динамическую аутентификацию транзакций; назван
по первым буквам Europay, Mastercard, Visa.
Спецификации поддерживаются
EMVCo. См. главу \ref{ch:emv}.

\section*{F}

\textbf{Force post} — отправка транзакции в клиринг
без онлайн-авторизации, с кодом авторизации, полученным
по телефону (voice
authorization). См. главу \ref{ch:tx-lifecycle}.

\textbf{FPAN} (Funding PAN) — настоящий PAN держателя карты,
лежащий в основе токена
(DPAN). См. главу \ref{ch:tokenization}.

\textbf{Friendly fraud} — вид мошенничества, при котором
легитимный держатель карты намеренно инициирует chargeback после
получения товара
или услуги. См. главу \ref{ch:disputes}.

\section*{H}

\textbf{Hold} — временная блокировка суммы на счёте держателя
карты, устанавливаемая эмитентом при одобрении авторизации;
снимается при settlement или по истечении
срока. См. главу \ref{ch:tx-lifecycle}.

\textbf{Hosted payment page} (HPP) — веб-страница,
размещённая на серверах платёжного шлюза, на которую мерчант
перенаправляет покупателя для ввода данных карты; мерчант
не видит, не принимает и не передаёт карточные
данные. См. главу \ref{ch:gateway}.

\textbf{HSM} (Hardware Security Module) — аппаратное
устройство, выполняющее криптографические операции (генерация
ключей, верификация PIN, проверка криптограмм) в защищённой,
физически изолированной среде. См. главу \ref{ch:crypto}.

\section*{I}

\textbf{iCVV} — код верификации, генерируемый EMV-чипом
с использованием другого сервисного кода; предотвращает
клонирование данных чипа на магнитную
полосу. См. главу \ref{ch:card-data}.

\textbf{IIN} (Issuer Identification Number) — первые
6 или 8 цифр PAN, идентифицирующие эмитента и карточную сеть;
формальное название по стандарту \iso{7812}, пришедшее на смену
термину BIN. См. главу \ref{ch:card-data}.

\textbf{Interchange} — межбанковская комиссия, которую
эквайер платит эмитенту за каждую карточную транзакцию;
рассчитывается на этапе клиринга и зависит от типа карты, MCC,
способа приёма и географии. См. главу \ref{ch:ecosystem}.

\textbf{IPM} (Integrated Processing Message) — формат
клиринговых и settlement-сообщений
\scheme{Mastercard}. См. главу \ref{ch:mastercard}.

\textbf{\iso{8583}} — международный стандарт формата сообщений
для обмена финансовой информацией между участниками карточного
платежа; определяет структуру MTI, bitmap и data
elements. См. главу \ref{ch:iso8583}.

\textbf{\iso{20022}} — стандарт финансовых сообщений
на основе XML/JSON, позиционируемый как преемник \iso{8583};
миграция идёт постепенно. См. главу \ref{ch:iso8583}.

\section*{K}

\textbf{KYC} (Know Your Customer) — процедура идентификации
и верификации клиента, обязательная для финансовых
организаций. См. главу \ref{ch:compliance-aml}.

\section*{L}

\textbf{Liability shift} — перенос финансовой ответственности
за мошенническую транзакцию на ту сторону (эмитент или эквайер),
которая не поддерживает EMV или 3-D Secure. См. главу \ref{ch:emv}.

\section*{M}

\textbf{MCC} (Merchant Category Code) — четырёхзначный код,
классифицирующий мерчанта по типу деятельности; влияет на ставку
interchange, сроки hold и правила
обработки. См. главу \ref{ch:tx-lifecycle}.

\textbf{MII} (Major Industry Identifier) — первая цифра PAN,
указывающая на отрасль эмитента (4 — \scheme{Visa},
5 — \scheme{Mastercard}, 2 — Мир). См. главу \ref{ch:card-data}.

\textbf{MSC} (Merchant Service Charge) — комиссия, которую
мерчант платит эквайеру за приём карт; включает interchange,
scheme fees и наценку эквайера. См. главу \ref{ch:ecosystem}.

\textbf{MTI} (Message Type Indicator) — первые четыре цифры
сообщения \iso{8583}, кодирующие версию стандарта, класс,
функцию и источник сообщения (например, \mti{0100} —
авторизационный запрос). См. главу \ref{ch:iso8583}.

\section*{N}

\textbf{NFC} (Near Field Communication) — технология
беспроводной связи ближнего действия, используемая
для бесконтактных платежей. См. главу \ref{ch:contactless}.

\section*{P}

\textbf{PAN} (Primary Account Number) — основной номер карты
(от 8 до 19 цифр), главный идентификатор карты во всей платёжной
инфраструктуре. Структура: IIN + номер счёта + контрольная цифра
(алгоритм Луна). См. главу \ref{ch:card-data}.

\textbf{Partial authorization} — одобрение транзакции
эмитентом не на полную запрошенную сумму, а на часть; покупатель
доплачивает оставшееся иным
способом. См. главу \ref{ch:tx-lifecycle}.

\textbf{Payment Facilitator} (PayFac) — организация,
зарегистрированная в карточной сети как мерчант-агрегатор,
процессящая транзакции sub-merchants под своим собственным MID
и берущая на себя onboarding, распределение средств
и compliance. См. главу \ref{ch:issuer-acquirer}.

\textbf{PCI DSS} (Payment Card Industry Data Security
Standard) — стандарт безопасности, определяющий требования
к защите cardholder data и sensitive authentication
data. См. главу \ref{ch:pci-dss}.

\textbf{Pre-authorization} (pre-auth) — авторизация
на предварительную сумму, когда финальная сумма заранее
неизвестна (гостиницы, аренда
автомобилей). См. главу \ref{ch:tx-lifecycle}.

\textbf{PSP} (Payment Service Provider) — компания,
предоставляющая мерчантам услуги по приёму платежей: API,
маршрутизацию, обработку авторизаций и settlement. На практике
PSP, платёжный шлюз и процессор часто совмещены в одном
продукте. См. главу \ref{ch:gateway}.

\section*{R}

\textbf{Reason code} — стандартизированный код, указывающий
причину chargeback или диспута; определяется правилами карточной
сети. См. главу \ref{ch:disputes}.

\textbf{Refund} — возврат денег держателю карты после
settlement; представляет собой новую транзакцию в обратном
направлении с полным циклом клиринга
и settlement. См. главу \ref{ch:tx-lifecycle}.

\textbf{Representment} — ответ мерчанта на chargeback
с предоставлением доказательств (compelling evidence)
легитимности транзакции. См. главу \ref{ch:disputes}.

\textbf{Response code} — двухсимвольный код ответа эмитента
в \de{39} сообщения \iso{8583} (например, \rc{00} — approved,
\rc{05} — do not honor, \rc{51} — insufficient
funds). См. главу \ref{ch:iso8583}.

\textbf{Reversal} — сообщение (MTI \mti{0400}), отменяющее
авторизацию до клиринга; мгновенно снимает hold со счёта
держателя. Самый \enquote{чистый} и быстрый способ отмены
платежа. См. главу \ref{ch:tx-lifecycle}.

\textbf{RRN} (Retrieval Reference Number) —
двенадцатисимвольный идентификатор транзакции, используемый
для поиска при запросах, диспутах
и сверке. См. главу \ref{ch:iso8583}.

\section*{S}

\textbf{SAD} (Sensitive Authentication Data) — чувствительные
данные аутентификации (полные данные Track 1/Track 2, CVV2, PIN
и PIN-блок); PCI DSS запрещает хранить их после авторизации
в любом виде. См. главу \ref{ch:card-data}.

\textbf{Scheme fees} — комиссии карточной сети (assessments):
плата за обработку транзакций, за трансграничную маршрутизацию
и за использование специальных
сервисов. См. главу \ref{ch:tx-lifecycle}.

\textbf{SDA} (Static Data Authentication) — базовый метод
офлайн-аутентификации EMV-чипа, проверяющий неизменность
подписанных данных с момента
персонализации. См. главу \ref{ch:emv}.

\textbf{Settlement} — фактическое перемещение денег между
банками по итогам клиринга; выполняется через нетто-расчёт
с участием расчётного
банка. См. главу \ref{ch:tx-lifecycle}.

\textbf{Settlement currency} — валюта, в которой карточная
сеть рассчитывается с эквайером; фиксируется в \de{50}
сообщения \iso{8583}. См. главу \ref{ch:multicurrency}.

\textbf{Single message} — архитектура обработки транзакций,
в которой авторизация и финансовая запись совмещены в одном
сообщении (MTI \mti{0200}/\mti{0210}); стандарт для банкоматов
и PIN-дебетовых сетей. См. главу \ref{ch:tx-lifecycle}.

\textbf{STAN} (Systems Trace Audit Number) — шестизначный
порядковый номер, присваиваемый каждой транзакции для привязки
запроса к ответу; уникален в пределах дня
и терминала. См. главу \ref{ch:iso8583}.

\section*{T}

\textbf{TADG} (Terminal and Application Design Guide) —
технический документ \scheme{Visa}, описывающий требования
к интерфейсу терминала, порядку выбора CVM и поведению
при различных статусах
транзакции. См. главу \ref{ch:visa}.

\textbf{TC} (Transaction Certificate) — криптограмма EMV,
означающая \enquote{карта одобряет транзакцию офлайн}; эмитент
не участвует в решении. См. главу \ref{ch:emv}.

\textbf{3-D Secure} (3DS) — протокол аутентификации держателя
карты в CNP-операциях, основанный на трёхдоменной модели: домен
эмитента (ACS), домен эквайера (3DS Server) и домен
совместимости (DS). См. главу \ref{ch:3ds}.

\textbf{TLV} (Tag-Length-Value) — формат кодирования данных
в \de{55} сообщения \iso{8583}: каждый объект содержит тег
(идентификатор), длину и значение. Используется для передачи
EMV-данных между чипом
и эмитентом. См. главу \ref{ch:iso8583}.

\textbf{TMS} (Terminal Management System) — система
удалённого управления парком POS-терминалов: загрузка ПО,
инъекция криптографических ключей, настройка параметров
и мониторинг. См. главу \ref{ch:issuer-acquirer}.

\textbf{Token vault} — защищённое хранилище, связывающее
токены (DPAN) с оригинальными PAN; управляется Token Service
Provider. См. главу \ref{ch:tokenization}.

\textbf{Track 1} (IATA) — первая дорожка магнитной полосы;
единственная, хранящая буквенно-цифровые данные, включая имя
держателя; до 79 символов. См. главу \ref{ch:card-data}.

\textbf{Track 2} (ABA) — вторая дорожка магнитной полосы,
основная рабочая дорожка для карточных платежей; только цифровые
данные (PAN, срок действия, сервисный код, CVV);
до 40 символов. См. главу \ref{ch:card-data}.

\textbf{Transaction currency} — валюта, в которой мерчант
выставляет цену; фиксируется в \de{49} сообщения
\iso{8583}. См. главу \ref{ch:multicurrency}.

\textbf{TSP} (Token Service Provider) — организация (обычно
карточная сеть), управляющая выпуском токенов, ведением token
vault и де-токенизацией. См. главу \ref{ch:tokenization}.

\textbf{TVR} (Terminal Verification Results) — пятибайтовое
поле EMV, фиксирующее результаты всех проверок, выполненных
терминалом (офлайн-аутентификация, верификация держателя,
риск-менеджмент). См. главу \ref{ch:emv}.

\section*{V}

\textbf{VCR} (Visa Claims Resolution) — система управления
диспутами \scheme{Visa} со структурированными reason codes
и ускоренными
сроками. См. главу \ref{ch:disputes}.

\textbf{Void} — аннулирование транзакции после её попадания
в клиринг, но до settlement; медленнее reversal, но быстрее
refund. См. главу \ref{ch:tx-lifecycle}.

%% ═══════════════════════════════════════════════════════════════
%%  КИРИЛЛИЦА  А – Я
%% ═══════════════════════════════════════════════════════════════

\section*{А}

\textbf{Авторизация} — фаза карточного платежа, происходящая
в реальном времени: эмитент проверяет карту, баланс, лимиты
и антифрод-правила, принимает решение об одобрении или отклонении
и устанавливает hold на счёте
держателя. См. главу \ref{ch:tx-lifecycle}.

\textbf{Алгоритм Луна} — метод вычисления контрольной цифры
PAN для обнаружения случайных ошибок ввода; не является
криптографической защитой. Запатентован Хансом Петером Луном
(IBM) в 1954 году. См. главу \ref{ch:card-data}.

\section*{Б}

\textbf{Billing currency} — валюта, в которой эмитент
списывает средства со счёта держателя карты; определяется
валютой карточного счёта. См. главу \ref{ch:multicurrency}.

\section*{К}

\textbf{Клиринг} — пакетная обработка транзакций: эквайер
формирует клиринговые файлы, карточная сеть сопоставляет их
с авторизациями и рассчитывает interchange
и scheme fees. См. главу \ref{ch:tx-lifecycle}.

\textbf{Корреспондентский счёт} — счёт, который один банк
открывает в другом для межбанковских расчётов. Ностро-счёт
(\emph{nostro} — \enquote{наш}) — наш счёт в чужом банке;
лоро-счёт (\emph{loro} — \enquote{их}) — тот же счёт
с точки зрения банка, в котором он
открыт. См. главу \ref{ch:banks}.

\section*{Н}

\textbf{НСПК} (АО \enquote{Национальная система платёжных
  карт}) — организация, выполняющая функции операционного
и платёжного клирингового центра (ОПКЦ) для внутрироссийских
транзакций по всем карточным системам, а также оператор
платёжной системы \enquote{Мир}. См. главу \ref{ch:mir}.

\section*{П}

\textbf{Платёжный оркестратор} (payment orchestrator) —
надстройка, маршрутизирующая транзакции между несколькими
шлюзами или процессорами по заданным
правилам. См. главу \ref{ch:gateway}.

\textbf{Платёжный процессор} (payment processor) — система,
непосредственно обрабатывающая авторизацию, клиринг
и settlement, часто имеющая прямое подключение к карточным
сетям. См. главу \ref{ch:gateway}.

\textbf{Платёжный шлюз} (payment gateway) — система,
принимающая платёжные запросы от мерчанта по высокоуровневому
API и транслирующая их в протоколы карточных сетей
и банков-эквайеров. См. главу \ref{ch:gateway}.

\textbf{ПЭП} (политически значимое лицо) — физическое лицо,
занимающее или занимавшее значительные публичные должности,
а также близкие родственники и известные соратники; для ПЭП
применяется Enhanced Due Diligence
(EDD). См. главу \ref{ch:compliance-aml}.

\textbf{Процессинговый центр} (processing center) — комплекс
программно-аппаратных средств, выполняющий авторизацию карточных
транзакций, формирование клиринговых файлов, расчёт
settlement-позиций и криптографические
операции. См. главу \ref{ch:processing-center}.

\section*{С}

\textbf{СБП} (Система быстрых платежей) — платёжная система
Банка России для мгновенных переводов между физическими лицами
и оплаты по QR-коду. См. главу \ref{ch:sbp}.

\textbf{Сервисный код} — три цифры на Track 1 и Track 2,
кодирующие тип обмена данными, тип авторизации и требования
к верификации; участвует в вычислении CVV
и iCVV. См. главу \ref{ch:card-data}.

\section*{Э}

\textbf{Эквайер} (acquirer) — банк или лицензированная
организация, обеспечивающая приём платёжных карт для мерчанта:
предоставляющая терминалы, маршрутизирующая авторизационные
запросы в карточную сеть и зачисляющая мерчанту средства
по итогам расчёта. См. главу \ref{ch:ecosystem}.

\textbf{Эмитент} (issuer) — банк или кредитная организация,
выпускающая платёжную карту и ведущая счёт, к которому она
привязана; принимает решение об одобрении или отклонении каждой
транзакции. См. главу \ref{ch:ecosystem}.
